%======================================================================
\subsection{BW-ROB-05 --- Motion study plan: Drive Joint, Motion Study, Motion Link (27 Oct 2026)}\label{sec:motion-plan}
%======================================================================

You chose \emph{kinematic first} on 27 Oct (\S\ref{sec:session-27oct}). This section is the plan the
\emph{motion} agent builds to; its Built subsection follows when it lands. It stays inside the Assembly track's
decisions of 19--25 Oct (\S\ref{sec:joints}: free drag, limits from datums, keep parts in place), and it reuses
what FreeCAD 1.1 already solves rather than writing a kinematics engine.

\paragraph{What exists, and what Fusion has on top.}
\begin{center}
\small
\begin{tabularx}{\linewidth}{@{}l>{\hsize=1.05\hsize}X>{\hsize=0.95\hsize}X@{}}
\toprule
Fusion & FreeCAD 1.1.4 (upstream) & Keystone after this epic\\
\midrule
Joint types (7) & 13 types incl.\ Gears, Belt, Rack \& pinion, Screw & all 13 on the ribbon, in one \emph{Joint} menu\\
Drive Joints & none; \code{Assembly\_CreateSimulation} needs a formula & \emph{Drive Joint}: slider over the joint's range, live\\
Motion Study & Simulation object, one Motion per joint, \code{formula(time)}, a scrub slider & \emph{Motion Study}: keyframe table, play/loop/ping-pong, speed, frames to PNG/MP4\\
Motion Link & Gears / RackPinion / Belt / Screw joints by radii or pitch & \emph{Motion Link}: pick two joints and a ratio; the right stock joint is made\\
Contact sets & none & \eref{ASM-10} (dynamics), later\\
\bottomrule
\end{tabularx}
\end{center}

\paragraph{Drive Joint.} A task panel in the right panel: the joint (a combo of the active assembly's movable
joints: Revolute, Slider, Cylindrical, Screw; Ball has no scalar coordinate), a slider with a spin box over the
joint's range (its limits when set, else $\pm180^\circ$ or $\pm100$\,mm), the model following live, and an
\emph{Animate} button sweeping the range back and forth at a chosen speed. \emph{OK} keeps the final pose as one
undo step; \emph{Cancel} restores the start pose and leaves no object in the document. The drive runs through the
same solver call the drag uses (the Python \code{preDrag}/\code{doDragStep}/\code{postDrag} from \S\ref{sec:joints}),
so limits and the drag backstop hold while driving.

\paragraph{Motion Study.} Fusion's dialog, as a task panel (figure~\ref{fig:motion-mock}): a row per joint,
keyframes along a time axis, the current frame highlighted, and the transport (play, pause, stop, loop,
ping-pong, speed $0.25\times$--$4\times$, a scrub slider). It is persisted as a Simulation object with one Motion per
joint, so the file opens in plain FreeCAD too; the keyframes themselves are a \code{[BW]} property on the study.
Each joint's motion is the piecewise-linear interpolation of its keyframes; the agent first tests whether the
solver's formula parser takes a piecewise expression, and otherwise steps the joints itself per frame with the
drive API. \emph{Save frames} renders the study off screen to a PNG sequence and, when \code{ffmpeg} is on the
path, an MP4.

\begin{figure}[htbp]
\centering
\begin{tikzpicture}[x=1mm,y=1mm,font=\sffamily\scriptsize]
  \draw[rounded corners=2mm, fill=white, draw=black!25] (0,0) rectangle (120,62);
  \node[anchor=west, font=\sffamily\small\bfseries] at (3,57.5) {Motion Study};
  \node[anchor=east, text=black!55] at (117,57.5) {2-DOF arm \quad duration 2.0 s};
  \draw[black!15] (0,54) -- (120,54);
  % transport
  \foreach \x/\t in {3/$\blacktriangleright$,11/$\blacksquare$,19/$\circlearrowleft$,27/$\leftrightarrow$}{
    \draw[rounded corners=0.8mm, draw=black!30] (\x,47.5) rectangle (\x+7,52.5); \node at (\x+3.5,50) {\t};}
  \node[anchor=west, text=black!55] at (37,50) {speed};
  \draw[rounded corners=0.8mm, draw=black!30] (46,47.5) rectangle (58,52.5); \node at (52,50) {1.0$\times$};
  \node[anchor=west, text=black!55] at (62,50) {t = 0.75 s};
  \draw[rounded corners=1mm, fill=bwblue, draw=bwblue] (96,47) rectangle (117,53);
  \node[text=white, font=\sffamily\scriptsize\bfseries] at (106.5,50) {Save frames};
  % time axis
  \draw[black!15] (0,44) -- (120,44);
  \node[anchor=west, text=black!55] at (3,41) {joint};
  \foreach \x/\t in {38/0,58/0.5,78/1.0,98/1.5,118/2.0}{\node[text=black!55] at (\x,41) {\t};}
  \draw[black!20] (38,38.5) -- (118,38.5);
  \foreach \x in {38,58,78,98,118}{\draw[black!20] (\x,37.5) -- (\x,39.5);}
  % rows
  \node[anchor=west] at (3,32) {Shoulder (Revolute)}; \node[anchor=west, text=black!55] at (3,28.5) {$0^\circ \to 90^\circ$};
  \draw[black!35] (38,31) -- (118,31);
  \foreach \x in {38,78,118}{\fill[bwblue] (\x,31) circle (1.2);}
  \node[anchor=south, text=bwblue] at (38,32) {0$^\circ$}; \node[anchor=south, text=bwblue] at (78,32) {90$^\circ$}; \node[anchor=south, text=bwblue] at (118,32) {0$^\circ$};
  \node[anchor=west] at (3,20) {Elbow (Revolute)}; \node[anchor=west, text=black!55] at (3,16.5) {$0^\circ \to -45^\circ$};
  \draw[black!35] (38,19) -- (118,19);
  \foreach \x in {38,98,118}{\fill[bwblue] (\x,19) circle (1.2);}
  \node[anchor=south, text=bwblue] at (38,20) {0$^\circ$}; \node[anchor=south, text=bwblue] at (98,20) {$-45^\circ$}; \node[anchor=south, text=bwblue] at (118,20) {0$^\circ$};
  \node[anchor=west] at (3,10) {Gripper (Slider)}; \node[anchor=west, text=black!55] at (3,6.5) {0 $\to$ 12 mm};
  \draw[black!35] (38,9) -- (118,9);
  \foreach \x in {38,118}{\fill[bwblue] (\x,9) circle (1.2);}
  % playhead
  \draw[bwred, thick] (68,36) -- (68,4); \fill[bwred] (68,36) -- (66.5,38.5) -- (69.5,38.5) -- cycle;
  \node[anchor=west, text=black!45] at (3,2.5) {double-click a row to add a keyframe at the playhead; drag a dot to move it};
\end{tikzpicture}
\caption{The Motion Study panel as planned: a row per joint, keyframes as dots on a time axis, the playhead at
0.75\,s, the transport with loop and ping-pong, and frame export. Values between keyframes are interpolated
(linear by default, see Q2).}
\label{fig:motion-mock}
\end{figure}

\paragraph{Motion Link.} Pick two joints and a ratio. The command makes the stock joint the solver already
understands, so nothing new is asked of the solver:
\begin{center}
\small
\begin{tabular}{@{}lll@{}}
\toprule
Joints & Stock joint made & Ratio means\\
\midrule
Revolute + Revolute & Gears (radii $r_1:r_2$) & $\theta_2=-(r_1/r_2)\,\theta_1$ (\emph{Reverse} flips the sign)\\
Revolute + Slider & Rack and pinion (pitch radius $r$) & $x=r\,\theta$, given as mm per revolution $=2\pi r$\\
Revolute + Revolute, same axis sense & Belt (radii) & $\theta_2=(r_1/r_2)\,\theta_1$\\
Revolute + Slider, same axis & Screw (pitch $p$) & $x=p\,\theta/2\pi$\\
Slider + Slider & not in the solver & reported as unsupported (Q4)\\
\bottomrule
\end{tabular}
\end{center}
The joint's label says what it is: ``Link: Shoulder $\to$ Wrist (1:3)''.

\paragraph{Theory: driven kinematics.} The assembly's parts have poses $\mathbf q$ (6 coordinates each), the
joints impose $\Phi(\mathbf q)=\mathbf 0$. A study adds one \emph{driver} per keyframed joint,
$\Phi_d(\mathbf q,t)=\theta_j(\mathbf q)-f_j(t)=0$, so at each output time the solver finds $\mathbf q$ from
$[\Phi;\Phi_d]=\mathbf 0$ by Newton iteration from the previous frame (this is what the Ondsel kinematic run
does with a Motion's formula). The keyframe function $f_j$ is piecewise linear between keyframes
$(t_k,v_k)$: $f_j(t)=v_k+(v_{k+1}-v_k)\,\frac{t-t_k}{t_{k+1}-t_k}$. Linear motion has velocity jumps at the
keyframes; a smooth option uses the cubic Hermite segment with zero end velocities (``ease in/out''),
$f_j(t)=v_k+(v_{k+1}-v_k)\,(3s^2-2s^3)$, $s=(t-t_k)/(t_{k+1}-t_k)$, whose velocity is zero at every keyframe and
peaks at $1.5\times$ the linear one. The drag backstop of \S\ref{sec:joints} still applies: a keyframe past a limit
is clamped and shown in orange. A Motion Link adds the algebraic constraint $\theta_2+(r_1/r_2)\theta_1=0$ (or
$x-r\theta=0$) and takes one degree of freedom away, so driving one joint moves both; driving both with
inconsistent keyframes over-constrains the study and the pre-check says which pair conflicts.

\paragraph{Verification.} A planar 2-DOF arm (links $\ell_1,\ell_2$) built in code: forward kinematics
$x=\ell_1\cos\theta_1+\ell_2\cos(\theta_1+\theta_2)$, $y=\ell_1\sin\theta_1+\ell_2\sin(\theta_1+\theta_2)$ at
$t=0,\,0.5,\,1$\,s must match the interpolated keyframes within $10^{-3}$\,mm; a 1:3 motion link turns the second
joint $3\times$ the first (sign by \emph{Reverse}). GUI tests open the three panels, move a part with the slider,
and check that \emph{Cancel} leaves no object.

\paragraph{Questions for you.}
\begin{enumerate}[label=M\arabic*.]
  \item \textbf{One study or many.} One Motion Study per assembly (Fusion's way), or several named studies kept
        in the file (a ``Studies'' group in the Browser)?\\
        \emph{Recommendation:} several; the first one is created without asking.
  \item \textbf{Interpolation.} Linear between keyframes (Fusion) or ease in/out by default?\\
        \emph{Recommendation:} linear by default, ease in/out as a per-study switch.
  \item \textbf{Export.} PNG sequence always, MP4 when \code{ffmpeg} is found; or bundle \code{ffmpeg} in Keystone
        (LGPL build, about 30\,MB)?\\
        \emph{Recommendation:} find it on the path for now; bundle it with \eref{FND-02}'s release work.
  \item \textbf{Slider-to-slider links.} The solver has no such joint. Leave it out, or emulate it with a hidden
        revolute and two rack-and-pinion joints?\\
        \emph{Recommendation:} leave it out until a design needs it.
  \item \textbf{Where Drive Joint lives.} Also on the joint's context menu in the Browser and on double-click of a
        joint marker in the 3D view?\\
        \emph{Recommendation:} yes to the context menu; double-click keeps opening the joint's own panel.
\end{enumerate}
